\label{der:4-3}
设$A$的行空间为$U$。$A\vx=0$逐行就是“$\vx$与每一行正交”，也等价于与全部行的线性组合正交，所以$\ker A=U^\perp$。对任意$U\subseteq\R^n$，取其单位正交基并扩充为全空间单位正交基；前面的基生成$U$，其余基恰生成$U^\perp$，再取正交补回到$U$，即$(U^\perp)^\perp=U$。

因此，若$A,B$行空间相同，其核相同，齐次组同解；若齐次组同解，即两核相同，则对两边取正交补，得到两行空间相同。非齐次时还需右端配合，例如$x+y=1$与$x+y=2$的系数行空间相同而解直线不同，这说明增加的条件不能删除。
